% =========================================================
% ppt/src/minggu-09.tex
% Slide pertemuan 9
% Bab X — Instruksi Transfer Data dan Aritmatika
%
% Sub-CPMK 3.1 | CPMK03 | CPL03
% =========================================================

\documentclass[aspectratio=169,11pt]{beamer}
% =========================================================
% ppt/src/preamble-ppt.tex
% Tema, warna, paket, dan makro bersama seluruh deck slide.
%
% Dipakai oleh setiap berkas minggu-NN.tex dengan pola:
%   \documentclass[aspectratio=169,11pt]{beamer}
%   % =========================================================
% ppt/src/preamble-ppt.tex
% Tema, warna, paket, dan makro bersama seluruh deck slide.
%
% Dipakai oleh setiap berkas minggu-NN.tex dengan pola:
%   \documentclass[aspectratio=169,11pt]{beamer}
%   % =========================================================
% ppt/src/preamble-ppt.tex
% Tema, warna, paket, dan makro bersama seluruh deck slide.
%
% Dipakai oleh setiap berkas minggu-NN.tex dengan pola:
%   \documentclass[aspectratio=169,11pt]{beamer}
%   \input{preamble-ppt}
%
% Palet warna disalin dari buku_ajar/config/styles.tex agar
% identitas visual slide sama dengan buku ajar.
% =========================================================

% ---------------------------------------------------------
% 1. Paket
% ---------------------------------------------------------
\usepackage[T1]{fontenc}
\usepackage[utf8]{inputenc}
\usepackage[indonesian]{babel}
\usepackage{lmodern}
\usepackage{tikz}
\usetikzlibrary{arrows.meta, positioning, calc, shapes.geometric}
\usepackage{booktabs}
\usepackage{array}
\usepackage{amsmath}
\usepackage{amsfonts}
\usepackage{pifont}          % \ding untuk tanda checklist
\usepackage{listings}

% ---------------------------------------------------------
% 2. Palet warna OBE (identik dengan buku)
% ---------------------------------------------------------
\definecolor{obeBlue}{HTML}{1F3864}
\definecolor{obeBlueLight}{HTML}{DCE4F5}
\definecolor{obeTeal}{HTML}{0F6E6E}
\definecolor{obeTealLight}{HTML}{D9F0EE}
\definecolor{obeGreen}{HTML}{2E6B34}
\definecolor{obeGreenLight}{HTML}{DEEFDF}
\definecolor{obeAmber}{HTML}{8A5A00}
\definecolor{obeAmberLight}{HTML}{FDF1D6}
\definecolor{obeMaroon}{HTML}{7A2233}
\definecolor{obeMaroonLight}{HTML}{F7E1E5}
\definecolor{obeGray}{HTML}{4A4A4A}
\definecolor{obeGrayLight}{HTML}{F0F0F0}
\definecolor{codeBack}{HTML}{F7F7F4}
\definecolor{codeFrame}{HTML}{C9C9C0}

% ---------------------------------------------------------
% 3. Identitas mata kuliah (dipakai seluruh deck)
% ---------------------------------------------------------
\newcommand{\NamaMK}{Bahasa Pemrograman Tingkat Rendah}
\newcommand{\NamaMKShort}{Bhs.\ Pemrograman Tingkat Rendah}
\newcommand{\KodeMK}{MWCA010}
\newcommand{\Prodi}{Program Studi Teknik Informatika}
\newcommand{\Fakultas}{Fakultas Teknologi Informasi}
\newcommand{\PerguruanTinggi}{Universitas Bale Bandung}
\newcommand{\DosenPengampu}{Cecep Suwanda, S.Si., M.Kom.}
\newcommand{\SemesterTeks}{Semester Ganjil}

% \PertemuanTeks dan \JudulBab ditetapkan di masing-masing
% berkas minggu-NN.tex sebelum \begin{document}.
\providecommand{\PertemuanTeks}{Pertemuan}

% ---------------------------------------------------------
% 4. Tema beamer
% ---------------------------------------------------------
\usetheme{default}
\setbeamertemplate{navigation symbols}{}

\setbeamercolor{normal text}{fg=black!88}
\setbeamercolor{structure}{fg=obeBlue}
\setbeamercolor{alerted text}{fg=obeMaroon}

% --- judul slide: pita biru penuh lebar ---
\setbeamercolor{frametitle}{fg=white,bg=obeBlue}
\setbeamerfont{frametitle}{series=\bfseries,size=\large}
\setbeamertemplate{frametitle}{%
    \nointerlineskip%
    \begin{beamercolorbox}[%
        wd=\paperwidth,ht=2.9ex,dp=1.5ex,
        leftskip=0.4cm,rightskip=0.4cm]{frametitle}%
        \usebeamerfont{frametitle}\insertframetitle%
    \end{beamercolorbox}%
    \nointerlineskip%
    \hbox{\textcolor{obeTeal}{\rule{\paperwidth}{1.1pt}}}%
}

% --- kaki slide: identitas mata kuliah, pertemuan, nomor ---
\setbeamercolor{footline}{fg=white,bg=obeGray}
\setbeamerfont{footline}{size=\scriptsize,family=\sffamily}
\setbeamertemplate{footline}{%
    \leavevmode%
    \hbox{%
        \begin{beamercolorbox}[wd=0.42\paperwidth,ht=2.1ex,dp=1.1ex,
            leftskip=0.4cm]{footline}%
            \usebeamerfont{footline}%
            \KodeMK\ \textbullet\ \NamaMKShort%
        \end{beamercolorbox}%
        \begin{beamercolorbox}[wd=0.32\paperwidth,ht=2.1ex,dp=1.1ex,
            center]{footline}%
            \usebeamerfont{footline}\PertemuanTeks%
        \end{beamercolorbox}%
        \begin{beamercolorbox}[wd=0.26\paperwidth,ht=2.1ex,dp=1.1ex,
            rightskip=0.4cm]{footline}%
            \usebeamerfont{footline}%
            \insertframenumber\,/\,\inserttotalframenumber%
        \end{beamercolorbox}%
    }%
    \vskip0pt%
}

% --- daftar ---
\setbeamercolor{itemize item}{fg=obeTeal}
\setbeamercolor{itemize subitem}{fg=obeTeal}
\setbeamercolor{enumerate item}{fg=obeBlue}
\setbeamertemplate{itemize item}{\raise0.2ex\hbox{\tiny$\blacksquare$}}
\setbeamertemplate{itemize subitem}{\raise0.2ex\hbox{\scriptsize$\bullet$}}

% --- kotak (pengganti tcolorbox pada buku) ---
\setbeamercolor{block title}{bg=obeBlueLight,fg=obeBlue}
\setbeamercolor{block body}{bg=obeBlueLight!45,fg=black!88}
\setbeamercolor{block title alerted}{bg=obeMaroonLight,fg=obeMaroon}
\setbeamercolor{block body alerted}{bg=obeMaroonLight!45,fg=black!88}
\setbeamercolor{block title example}{bg=obeTealLight,fg=obeTeal}
\setbeamercolor{block body example}{bg=obeTealLight!45,fg=black!88}

% ---------------------------------------------------------
% 5. Makro notasi (sama persis dengan buku)
%
% Catatan: \note TIDAK didefinisikan di sini karena beamer
% sudah memakai \note untuk catatan pembicara. Untuk catatan
% di slide dipakai \catatan.
% ---------------------------------------------------------
\newcommand{\code}[1]{\texttt{#1}}
\newcommand{\asm}[1]{\texttt{#1}}
\newcommand{\reg}[1]{\texttt{#1}}
\newcommand{\flag}[1]{\texttt{#1}}
\newcommand{\dir}[1]{\texttt{#1}}
\newcommand{\file}[1]{\texttt{#1}}
\newcommand{\eng}[1]{\textit{#1}}

\newcommand{\hex}[1]{\texttt{#1h}}
\newcommand{\bin}[1]{\texttt{#1b}}
\newcommand{\dec}[1]{\texttt{#1}}

\newcommand{\important}[1]{\textbf{Penting:} #1}
\newcommand{\catatan}[1]{\textbf{Catatan:} #1}
\newcommand{\term}[1]{\textbf{#1}}

\newcommand{\CPL}[1]{\textbf{CPL#1}}
\newcommand{\CPMK}[1]{\textbf{CPMK#1}}
\newcommand{\SubCPMK}[1]{\textbf{Sub-CPMK #1}}

\newcommand{\checkbox}{\ding{113}\hspace{0.35em}}
\newcommand{\checked}{\ding{51}\hspace{0.35em}}
\newcommand{\crossed}{\ding{55}\hspace{0.35em}}

\newcommand{\flowarrow}{\ensuremath{\rightarrow}}

% ---------------------------------------------------------
% 6. Makro tata letak slide
% ---------------------------------------------------------

% Keterangan sumber materi (ketertelusuran ke buku)
\newcommand{\Sumber}[1]{%
    \vfill
    {\scriptsize\color{obeGray}Sumber: #1\par}%
}

% Frame judul deck: \JudulDeck{<judul bab>}{<pertemuan>}
\newcommand{\JudulDeck}[2]{%
\begin{frame}[plain]
    \vspace{0.25cm}
    {\color{obeBlue}\sffamily\bfseries\small \NamaMK\ (\KodeMK)\par}
    \vspace{0.12cm}
    {\color{obeGray}\sffamily\scriptsize
     \Prodi\ \textbullet\ \Fakultas\ \textbullet\ \PerguruanTinggi\par}
    \vfill
    {\color{obeBlue}\sffamily\LARGE\bfseries #1\par}
    \vspace{0.4cm}
    {\color{obeTeal}\sffamily\large #2\par}
    \vspace{0.45cm}
    {\color{obeTeal}\rule{\textwidth}{1pt}}
    \vfill
    {\color{obeGray}\sffamily\small Dosen Pengampu\par}
    \vspace{0.08cm}
    {\sffamily\bfseries \DosenPengampu\par}
    \vspace{0.2cm}
    {\color{obeGray}\sffamily\small \SemesterTeks\par}
    \vspace{0.25cm}
\end{frame}
}

% Frame pemisah bagian
\newcommand{\PemisahBagian}[1]{%
\begin{frame}[plain]
    \vfill
    \begin{center}
        {\color{obeTeal}\sffamily\small\bfseries BAGIAN\par}
        \vspace{0.35cm}
        {\color{obeBlue}\sffamily\Large\bfseries #1\par}
        \vspace{0.35cm}
        {\color{obeTeal}\rule{0.55\textwidth}{0.9pt}}
    \end{center}
    \vfill
\end{frame}
}

% ---------------------------------------------------------
% 7. Gaya penampil kode
% ---------------------------------------------------------
\input{listing-ppt}

%
% Palet warna disalin dari buku_ajar/config/styles.tex agar
% identitas visual slide sama dengan buku ajar.
% =========================================================

% ---------------------------------------------------------
% 1. Paket
% ---------------------------------------------------------
\usepackage[T1]{fontenc}
\usepackage[utf8]{inputenc}
\usepackage[indonesian]{babel}
\usepackage{lmodern}
\usepackage{tikz}
\usetikzlibrary{arrows.meta, positioning, calc, shapes.geometric}
\usepackage{booktabs}
\usepackage{array}
\usepackage{amsmath}
\usepackage{amsfonts}
\usepackage{pifont}          % \ding untuk tanda checklist
\usepackage{listings}

% ---------------------------------------------------------
% 2. Palet warna OBE (identik dengan buku)
% ---------------------------------------------------------
\definecolor{obeBlue}{HTML}{1F3864}
\definecolor{obeBlueLight}{HTML}{DCE4F5}
\definecolor{obeTeal}{HTML}{0F6E6E}
\definecolor{obeTealLight}{HTML}{D9F0EE}
\definecolor{obeGreen}{HTML}{2E6B34}
\definecolor{obeGreenLight}{HTML}{DEEFDF}
\definecolor{obeAmber}{HTML}{8A5A00}
\definecolor{obeAmberLight}{HTML}{FDF1D6}
\definecolor{obeMaroon}{HTML}{7A2233}
\definecolor{obeMaroonLight}{HTML}{F7E1E5}
\definecolor{obeGray}{HTML}{4A4A4A}
\definecolor{obeGrayLight}{HTML}{F0F0F0}
\definecolor{codeBack}{HTML}{F7F7F4}
\definecolor{codeFrame}{HTML}{C9C9C0}

% ---------------------------------------------------------
% 3. Identitas mata kuliah (dipakai seluruh deck)
% ---------------------------------------------------------
\newcommand{\NamaMK}{Bahasa Pemrograman Tingkat Rendah}
\newcommand{\NamaMKShort}{Bhs.\ Pemrograman Tingkat Rendah}
\newcommand{\KodeMK}{MWCA010}
\newcommand{\Prodi}{Program Studi Teknik Informatika}
\newcommand{\Fakultas}{Fakultas Teknologi Informasi}
\newcommand{\PerguruanTinggi}{Universitas Bale Bandung}
\newcommand{\DosenPengampu}{Cecep Suwanda, S.Si., M.Kom.}
\newcommand{\SemesterTeks}{Semester Ganjil}

% \PertemuanTeks dan \JudulBab ditetapkan di masing-masing
% berkas minggu-NN.tex sebelum \begin{document}.
\providecommand{\PertemuanTeks}{Pertemuan}

% ---------------------------------------------------------
% 4. Tema beamer
% ---------------------------------------------------------
\usetheme{default}
\setbeamertemplate{navigation symbols}{}

\setbeamercolor{normal text}{fg=black!88}
\setbeamercolor{structure}{fg=obeBlue}
\setbeamercolor{alerted text}{fg=obeMaroon}

% --- judul slide: pita biru penuh lebar ---
\setbeamercolor{frametitle}{fg=white,bg=obeBlue}
\setbeamerfont{frametitle}{series=\bfseries,size=\large}
\setbeamertemplate{frametitle}{%
    \nointerlineskip%
    \begin{beamercolorbox}[%
        wd=\paperwidth,ht=2.9ex,dp=1.5ex,
        leftskip=0.4cm,rightskip=0.4cm]{frametitle}%
        \usebeamerfont{frametitle}\insertframetitle%
    \end{beamercolorbox}%
    \nointerlineskip%
    \hbox{\textcolor{obeTeal}{\rule{\paperwidth}{1.1pt}}}%
}

% --- kaki slide: identitas mata kuliah, pertemuan, nomor ---
\setbeamercolor{footline}{fg=white,bg=obeGray}
\setbeamerfont{footline}{size=\scriptsize,family=\sffamily}
\setbeamertemplate{footline}{%
    \leavevmode%
    \hbox{%
        \begin{beamercolorbox}[wd=0.42\paperwidth,ht=2.1ex,dp=1.1ex,
            leftskip=0.4cm]{footline}%
            \usebeamerfont{footline}%
            \KodeMK\ \textbullet\ \NamaMKShort%
        \end{beamercolorbox}%
        \begin{beamercolorbox}[wd=0.32\paperwidth,ht=2.1ex,dp=1.1ex,
            center]{footline}%
            \usebeamerfont{footline}\PertemuanTeks%
        \end{beamercolorbox}%
        \begin{beamercolorbox}[wd=0.26\paperwidth,ht=2.1ex,dp=1.1ex,
            rightskip=0.4cm]{footline}%
            \usebeamerfont{footline}%
            \insertframenumber\,/\,\inserttotalframenumber%
        \end{beamercolorbox}%
    }%
    \vskip0pt%
}

% --- daftar ---
\setbeamercolor{itemize item}{fg=obeTeal}
\setbeamercolor{itemize subitem}{fg=obeTeal}
\setbeamercolor{enumerate item}{fg=obeBlue}
\setbeamertemplate{itemize item}{\raise0.2ex\hbox{\tiny$\blacksquare$}}
\setbeamertemplate{itemize subitem}{\raise0.2ex\hbox{\scriptsize$\bullet$}}

% --- kotak (pengganti tcolorbox pada buku) ---
\setbeamercolor{block title}{bg=obeBlueLight,fg=obeBlue}
\setbeamercolor{block body}{bg=obeBlueLight!45,fg=black!88}
\setbeamercolor{block title alerted}{bg=obeMaroonLight,fg=obeMaroon}
\setbeamercolor{block body alerted}{bg=obeMaroonLight!45,fg=black!88}
\setbeamercolor{block title example}{bg=obeTealLight,fg=obeTeal}
\setbeamercolor{block body example}{bg=obeTealLight!45,fg=black!88}

% ---------------------------------------------------------
% 5. Makro notasi (sama persis dengan buku)
%
% Catatan: \note TIDAK didefinisikan di sini karena beamer
% sudah memakai \note untuk catatan pembicara. Untuk catatan
% di slide dipakai \catatan.
% ---------------------------------------------------------
\newcommand{\code}[1]{\texttt{#1}}
\newcommand{\asm}[1]{\texttt{#1}}
\newcommand{\reg}[1]{\texttt{#1}}
\newcommand{\flag}[1]{\texttt{#1}}
\newcommand{\dir}[1]{\texttt{#1}}
\newcommand{\file}[1]{\texttt{#1}}
\newcommand{\eng}[1]{\textit{#1}}

\newcommand{\hex}[1]{\texttt{#1h}}
\newcommand{\bin}[1]{\texttt{#1b}}
\newcommand{\dec}[1]{\texttt{#1}}

\newcommand{\important}[1]{\textbf{Penting:} #1}
\newcommand{\catatan}[1]{\textbf{Catatan:} #1}
\newcommand{\term}[1]{\textbf{#1}}

\newcommand{\CPL}[1]{\textbf{CPL#1}}
\newcommand{\CPMK}[1]{\textbf{CPMK#1}}
\newcommand{\SubCPMK}[1]{\textbf{Sub-CPMK #1}}

\newcommand{\checkbox}{\ding{113}\hspace{0.35em}}
\newcommand{\checked}{\ding{51}\hspace{0.35em}}
\newcommand{\crossed}{\ding{55}\hspace{0.35em}}

\newcommand{\flowarrow}{\ensuremath{\rightarrow}}

% ---------------------------------------------------------
% 6. Makro tata letak slide
% ---------------------------------------------------------

% Keterangan sumber materi (ketertelusuran ke buku)
\newcommand{\Sumber}[1]{%
    \vfill
    {\scriptsize\color{obeGray}Sumber: #1\par}%
}

% Frame judul deck: \JudulDeck{<judul bab>}{<pertemuan>}
\newcommand{\JudulDeck}[2]{%
\begin{frame}[plain]
    \vspace{0.25cm}
    {\color{obeBlue}\sffamily\bfseries\small \NamaMK\ (\KodeMK)\par}
    \vspace{0.12cm}
    {\color{obeGray}\sffamily\scriptsize
     \Prodi\ \textbullet\ \Fakultas\ \textbullet\ \PerguruanTinggi\par}
    \vfill
    {\color{obeBlue}\sffamily\LARGE\bfseries #1\par}
    \vspace{0.4cm}
    {\color{obeTeal}\sffamily\large #2\par}
    \vspace{0.45cm}
    {\color{obeTeal}\rule{\textwidth}{1pt}}
    \vfill
    {\color{obeGray}\sffamily\small Dosen Pengampu\par}
    \vspace{0.08cm}
    {\sffamily\bfseries \DosenPengampu\par}
    \vspace{0.2cm}
    {\color{obeGray}\sffamily\small \SemesterTeks\par}
    \vspace{0.25cm}
\end{frame}
}

% Frame pemisah bagian
\newcommand{\PemisahBagian}[1]{%
\begin{frame}[plain]
    \vfill
    \begin{center}
        {\color{obeTeal}\sffamily\small\bfseries BAGIAN\par}
        \vspace{0.35cm}
        {\color{obeBlue}\sffamily\Large\bfseries #1\par}
        \vspace{0.35cm}
        {\color{obeTeal}\rule{0.55\textwidth}{0.9pt}}
    \end{center}
    \vfill
\end{frame}
}

% ---------------------------------------------------------
% 7. Gaya penampil kode
% ---------------------------------------------------------
% =========================================================
% ppt/src/listing-ppt.tex
% Gaya penampil kode untuk slide.
%
% Diturunkan dari buku_ajar/config/listings.tex, tetapi
% disesuaikan untuk slide: ukuran lebih kecil (karena frame
% 16:9 dan kode harus utuh terlihat) dan jarak lebih rapat.
% =========================================================

% ---------------------------------------------------------
% 1. Definisi bahasa MASM/TASM 8086 (salinan dari buku)
% ---------------------------------------------------------
\lstdefinelanguage{MASM}{
    morekeywords={
        MOV, XCHG, LEA, LDS, LES, PUSH, POP, PUSHA, POPA, PUSHF, POPF,
        LAHF, SAHF, XLAT, IN, OUT, CBW, CWD,
        ADD, ADC, SUB, SBB, INC, DEC, NEG, CMP, MUL, IMUL, DIV, IDIV,
        AAA, AAS, AAM, AAD, DAA, DAS,
        AND, OR, XOR, NOT, TEST, SHL, SHR, SAL, SAR, ROL, ROR, RCL, RCR,
        JMP, JE, JZ, JNE, JNZ, JG, JGE, JL, JLE, JA, JAE, JB, JBE,
        JC, JNC, JS, JNS, JO, JNO, JP, JPE, JNP, JPO, JCXZ,
        LOOP, LOOPE, LOOPZ, LOOPNE, LOOPNZ, CALL, RET, INT, IRET, INTO,
        HLT, NOP, CLI, STI, CLD, STD, CLC, STC, CMC, WAIT, ESC, LOCK,
        DB, DW, DD, DQ, DT, EQU, ORG, SEGMENT, ENDS, ASSUME, END,
        PROC, ENDP, MACRO, ENDM, LOCAL, PUBLIC, EXTRN, EXTERN,
        INCLUDE, INCLUDELIB, MODEL, STACK, CODE, DATA, CONST, ALIGN, EVEN,
        OFFSET, PTR, BYTE, WORD, DWORD, QWORD, TBYTE,
        NEAR, FAR, SHORT, LABEL, TYPE, SIZE, LENGTH, DUP, THIS,
        AX, BX, CX, DX, AL, AH, BL, BH, CL, CH, DL, DH,
        SP, BP, SI, DI, IP, CS, DS, SS, ES, FLAGS,
        CF, PF, AF, ZF, TF, IF, DF, OF
    },
    sensitive=false,
    morecomment=[l]{;},
    morestring=[b]',
    morestring=[b]"
}

% ---------------------------------------------------------
% 2. Gaya kode assembly untuk slide
% ---------------------------------------------------------
\lstdefinestyle{asmslide}{
    language=MASM,
    basicstyle=\ttfamily\scriptsize,
    keywordstyle=\color{obeBlue}\bfseries,
    commentstyle=\color{obeGreen}\itshape,
    stringstyle=\color{obeMaroon},
    numberstyle=\tiny\color{obeGray},
    numbers=left,
    numbersep=5pt,
    stepnumber=1,
    backgroundcolor=\color{codeBack},
    frame=single,
    rulecolor=\color{codeFrame},
    framesep=4pt,
    breaklines=true,
    breakatwhitespace=false,
    % CATATAN PENTING: opsi `postbreak` TIDAK boleh dipakai di sini.
    % Di beamer, `breaklines=true` bersama `postbreak=<kode>` membuat
    % listings gagal dirakit dengan galat
    %   "Missing number, treated as zero <to be read again> \kern".
    % Karena \lstset{style=asmslide} bersifat global, opsi postbreak
    % juga bocor ke gaya lain (mis. outslide/keluaran) dan merusaknya.
    % Di buku ajar (kelas `book`) justru sebaliknya: postbreak aman dan
    % dipakai. Jadi larangan ini khusus untuk slide.
    showstringspaces=false,
    tabsize=4,
    columns=fullflexible,
    keepspaces=true,
    upquote=true,
    xleftmargin=12pt,
    aboveskip=0.4em,
    belowskip=0.4em
}

% ---------------------------------------------------------
% 3. Gaya keluaran program / sesi debugger untuk slide
% ---------------------------------------------------------
\lstdefinestyle{outslide}{
    basicstyle=\ttfamily\scriptsize,
    keywordstyle=\color{black},
    commentstyle=\color{obeGray},
    stringstyle=\color{black},
    numbers=none,
    backgroundcolor=\color{obeGrayLight},
    frame=single,
    rulecolor=\color{codeFrame},
    framesep=4pt,
    breaklines=true,
    showstringspaces=false,
    columns=fullflexible,
    keepspaces=true,
    upquote=true,
    xleftmargin=4pt,
    aboveskip=0.4em,
    belowskip=0.4em
}

\lstset{style=asmslide}

% ---------------------------------------------------------
% 4. Environment ringkas untuk dipakai di dalam frame
%
% Frame yang memuat kedua environment ini WAJIB memakai
% opsi [fragile], karena isinya verbatim.
% ---------------------------------------------------------
\lstnewenvironment{kodeasm}[1][]{%
    \lstset{style=asmslide, #1}%
}{}

\lstnewenvironment{keluaran}[1][]{%
    \lstset{style=outslide, #1}%
}{}


%
% Palet warna disalin dari buku_ajar/config/styles.tex agar
% identitas visual slide sama dengan buku ajar.
% =========================================================

% ---------------------------------------------------------
% 1. Paket
% ---------------------------------------------------------
\usepackage[T1]{fontenc}
\usepackage[utf8]{inputenc}
\usepackage[indonesian]{babel}
\usepackage{lmodern}
\usepackage{tikz}
\usetikzlibrary{arrows.meta, positioning, calc, shapes.geometric}
\usepackage{booktabs}
\usepackage{array}
\usepackage{amsmath}
\usepackage{amsfonts}
\usepackage{pifont}          % \ding untuk tanda checklist
\usepackage{listings}

% ---------------------------------------------------------
% 2. Palet warna OBE (identik dengan buku)
% ---------------------------------------------------------
\definecolor{obeBlue}{HTML}{1F3864}
\definecolor{obeBlueLight}{HTML}{DCE4F5}
\definecolor{obeTeal}{HTML}{0F6E6E}
\definecolor{obeTealLight}{HTML}{D9F0EE}
\definecolor{obeGreen}{HTML}{2E6B34}
\definecolor{obeGreenLight}{HTML}{DEEFDF}
\definecolor{obeAmber}{HTML}{8A5A00}
\definecolor{obeAmberLight}{HTML}{FDF1D6}
\definecolor{obeMaroon}{HTML}{7A2233}
\definecolor{obeMaroonLight}{HTML}{F7E1E5}
\definecolor{obeGray}{HTML}{4A4A4A}
\definecolor{obeGrayLight}{HTML}{F0F0F0}
\definecolor{codeBack}{HTML}{F7F7F4}
\definecolor{codeFrame}{HTML}{C9C9C0}

% ---------------------------------------------------------
% 3. Identitas mata kuliah (dipakai seluruh deck)
% ---------------------------------------------------------
\newcommand{\NamaMK}{Bahasa Pemrograman Tingkat Rendah}
\newcommand{\NamaMKShort}{Bhs.\ Pemrograman Tingkat Rendah}
\newcommand{\KodeMK}{MWCA010}
\newcommand{\Prodi}{Program Studi Teknik Informatika}
\newcommand{\Fakultas}{Fakultas Teknologi Informasi}
\newcommand{\PerguruanTinggi}{Universitas Bale Bandung}
\newcommand{\DosenPengampu}{Cecep Suwanda, S.Si., M.Kom.}
\newcommand{\SemesterTeks}{Semester Ganjil}

% \PertemuanTeks dan \JudulBab ditetapkan di masing-masing
% berkas minggu-NN.tex sebelum \begin{document}.
\providecommand{\PertemuanTeks}{Pertemuan}

% ---------------------------------------------------------
% 4. Tema beamer
% ---------------------------------------------------------
\usetheme{default}
\setbeamertemplate{navigation symbols}{}

\setbeamercolor{normal text}{fg=black!88}
\setbeamercolor{structure}{fg=obeBlue}
\setbeamercolor{alerted text}{fg=obeMaroon}

% --- judul slide: pita biru penuh lebar ---
\setbeamercolor{frametitle}{fg=white,bg=obeBlue}
\setbeamerfont{frametitle}{series=\bfseries,size=\large}
\setbeamertemplate{frametitle}{%
    \nointerlineskip%
    \begin{beamercolorbox}[%
        wd=\paperwidth,ht=2.9ex,dp=1.5ex,
        leftskip=0.4cm,rightskip=0.4cm]{frametitle}%
        \usebeamerfont{frametitle}\insertframetitle%
    \end{beamercolorbox}%
    \nointerlineskip%
    \hbox{\textcolor{obeTeal}{\rule{\paperwidth}{1.1pt}}}%
}

% --- kaki slide: identitas mata kuliah, pertemuan, nomor ---
\setbeamercolor{footline}{fg=white,bg=obeGray}
\setbeamerfont{footline}{size=\scriptsize,family=\sffamily}
\setbeamertemplate{footline}{%
    \leavevmode%
    \hbox{%
        \begin{beamercolorbox}[wd=0.42\paperwidth,ht=2.1ex,dp=1.1ex,
            leftskip=0.4cm]{footline}%
            \usebeamerfont{footline}%
            \KodeMK\ \textbullet\ \NamaMKShort%
        \end{beamercolorbox}%
        \begin{beamercolorbox}[wd=0.32\paperwidth,ht=2.1ex,dp=1.1ex,
            center]{footline}%
            \usebeamerfont{footline}\PertemuanTeks%
        \end{beamercolorbox}%
        \begin{beamercolorbox}[wd=0.26\paperwidth,ht=2.1ex,dp=1.1ex,
            rightskip=0.4cm]{footline}%
            \usebeamerfont{footline}%
            \insertframenumber\,/\,\inserttotalframenumber%
        \end{beamercolorbox}%
    }%
    \vskip0pt%
}

% --- daftar ---
\setbeamercolor{itemize item}{fg=obeTeal}
\setbeamercolor{itemize subitem}{fg=obeTeal}
\setbeamercolor{enumerate item}{fg=obeBlue}
\setbeamertemplate{itemize item}{\raise0.2ex\hbox{\tiny$\blacksquare$}}
\setbeamertemplate{itemize subitem}{\raise0.2ex\hbox{\scriptsize$\bullet$}}

% --- kotak (pengganti tcolorbox pada buku) ---
\setbeamercolor{block title}{bg=obeBlueLight,fg=obeBlue}
\setbeamercolor{block body}{bg=obeBlueLight!45,fg=black!88}
\setbeamercolor{block title alerted}{bg=obeMaroonLight,fg=obeMaroon}
\setbeamercolor{block body alerted}{bg=obeMaroonLight!45,fg=black!88}
\setbeamercolor{block title example}{bg=obeTealLight,fg=obeTeal}
\setbeamercolor{block body example}{bg=obeTealLight!45,fg=black!88}

% ---------------------------------------------------------
% 5. Makro notasi (sama persis dengan buku)
%
% Catatan: \note TIDAK didefinisikan di sini karena beamer
% sudah memakai \note untuk catatan pembicara. Untuk catatan
% di slide dipakai \catatan.
% ---------------------------------------------------------
\newcommand{\code}[1]{\texttt{#1}}
\newcommand{\asm}[1]{\texttt{#1}}
\newcommand{\reg}[1]{\texttt{#1}}
\newcommand{\flag}[1]{\texttt{#1}}
\newcommand{\dir}[1]{\texttt{#1}}
\newcommand{\file}[1]{\texttt{#1}}
\newcommand{\eng}[1]{\textit{#1}}

\newcommand{\hex}[1]{\texttt{#1h}}
\newcommand{\bin}[1]{\texttt{#1b}}
\newcommand{\dec}[1]{\texttt{#1}}

\newcommand{\important}[1]{\textbf{Penting:} #1}
\newcommand{\catatan}[1]{\textbf{Catatan:} #1}
\newcommand{\term}[1]{\textbf{#1}}

\newcommand{\CPL}[1]{\textbf{CPL#1}}
\newcommand{\CPMK}[1]{\textbf{CPMK#1}}
\newcommand{\SubCPMK}[1]{\textbf{Sub-CPMK #1}}

\newcommand{\checkbox}{\ding{113}\hspace{0.35em}}
\newcommand{\checked}{\ding{51}\hspace{0.35em}}
\newcommand{\crossed}{\ding{55}\hspace{0.35em}}

\newcommand{\flowarrow}{\ensuremath{\rightarrow}}

% ---------------------------------------------------------
% 6. Makro tata letak slide
% ---------------------------------------------------------

% Keterangan sumber materi (ketertelusuran ke buku)
\newcommand{\Sumber}[1]{%
    \vfill
    {\scriptsize\color{obeGray}Sumber: #1\par}%
}

% Frame judul deck: \JudulDeck{<judul bab>}{<pertemuan>}
\newcommand{\JudulDeck}[2]{%
\begin{frame}[plain]
    \vspace{0.25cm}
    {\color{obeBlue}\sffamily\bfseries\small \NamaMK\ (\KodeMK)\par}
    \vspace{0.12cm}
    {\color{obeGray}\sffamily\scriptsize
     \Prodi\ \textbullet\ \Fakultas\ \textbullet\ \PerguruanTinggi\par}
    \vfill
    {\color{obeBlue}\sffamily\LARGE\bfseries #1\par}
    \vspace{0.4cm}
    {\color{obeTeal}\sffamily\large #2\par}
    \vspace{0.45cm}
    {\color{obeTeal}\rule{\textwidth}{1pt}}
    \vfill
    {\color{obeGray}\sffamily\small Dosen Pengampu\par}
    \vspace{0.08cm}
    {\sffamily\bfseries \DosenPengampu\par}
    \vspace{0.2cm}
    {\color{obeGray}\sffamily\small \SemesterTeks\par}
    \vspace{0.25cm}
\end{frame}
}

% Frame pemisah bagian
\newcommand{\PemisahBagian}[1]{%
\begin{frame}[plain]
    \vfill
    \begin{center}
        {\color{obeTeal}\sffamily\small\bfseries BAGIAN\par}
        \vspace{0.35cm}
        {\color{obeBlue}\sffamily\Large\bfseries #1\par}
        \vspace{0.35cm}
        {\color{obeTeal}\rule{0.55\textwidth}{0.9pt}}
    \end{center}
    \vfill
\end{frame}
}

% ---------------------------------------------------------
% 7. Gaya penampil kode
% ---------------------------------------------------------
% =========================================================
% ppt/src/listing-ppt.tex
% Gaya penampil kode untuk slide.
%
% Diturunkan dari buku_ajar/config/listings.tex, tetapi
% disesuaikan untuk slide: ukuran lebih kecil (karena frame
% 16:9 dan kode harus utuh terlihat) dan jarak lebih rapat.
% =========================================================

% ---------------------------------------------------------
% 1. Definisi bahasa MASM/TASM 8086 (salinan dari buku)
% ---------------------------------------------------------
\lstdefinelanguage{MASM}{
    morekeywords={
        MOV, XCHG, LEA, LDS, LES, PUSH, POP, PUSHA, POPA, PUSHF, POPF,
        LAHF, SAHF, XLAT, IN, OUT, CBW, CWD,
        ADD, ADC, SUB, SBB, INC, DEC, NEG, CMP, MUL, IMUL, DIV, IDIV,
        AAA, AAS, AAM, AAD, DAA, DAS,
        AND, OR, XOR, NOT, TEST, SHL, SHR, SAL, SAR, ROL, ROR, RCL, RCR,
        JMP, JE, JZ, JNE, JNZ, JG, JGE, JL, JLE, JA, JAE, JB, JBE,
        JC, JNC, JS, JNS, JO, JNO, JP, JPE, JNP, JPO, JCXZ,
        LOOP, LOOPE, LOOPZ, LOOPNE, LOOPNZ, CALL, RET, INT, IRET, INTO,
        HLT, NOP, CLI, STI, CLD, STD, CLC, STC, CMC, WAIT, ESC, LOCK,
        DB, DW, DD, DQ, DT, EQU, ORG, SEGMENT, ENDS, ASSUME, END,
        PROC, ENDP, MACRO, ENDM, LOCAL, PUBLIC, EXTRN, EXTERN,
        INCLUDE, INCLUDELIB, MODEL, STACK, CODE, DATA, CONST, ALIGN, EVEN,
        OFFSET, PTR, BYTE, WORD, DWORD, QWORD, TBYTE,
        NEAR, FAR, SHORT, LABEL, TYPE, SIZE, LENGTH, DUP, THIS,
        AX, BX, CX, DX, AL, AH, BL, BH, CL, CH, DL, DH,
        SP, BP, SI, DI, IP, CS, DS, SS, ES, FLAGS,
        CF, PF, AF, ZF, TF, IF, DF, OF
    },
    sensitive=false,
    morecomment=[l]{;},
    morestring=[b]',
    morestring=[b]"
}

% ---------------------------------------------------------
% 2. Gaya kode assembly untuk slide
% ---------------------------------------------------------
\lstdefinestyle{asmslide}{
    language=MASM,
    basicstyle=\ttfamily\scriptsize,
    keywordstyle=\color{obeBlue}\bfseries,
    commentstyle=\color{obeGreen}\itshape,
    stringstyle=\color{obeMaroon},
    numberstyle=\tiny\color{obeGray},
    numbers=left,
    numbersep=5pt,
    stepnumber=1,
    backgroundcolor=\color{codeBack},
    frame=single,
    rulecolor=\color{codeFrame},
    framesep=4pt,
    breaklines=true,
    breakatwhitespace=false,
    % CATATAN PENTING: opsi `postbreak` TIDAK boleh dipakai di sini.
    % Di beamer, `breaklines=true` bersama `postbreak=<kode>` membuat
    % listings gagal dirakit dengan galat
    %   "Missing number, treated as zero <to be read again> \kern".
    % Karena \lstset{style=asmslide} bersifat global, opsi postbreak
    % juga bocor ke gaya lain (mis. outslide/keluaran) dan merusaknya.
    % Di buku ajar (kelas `book`) justru sebaliknya: postbreak aman dan
    % dipakai. Jadi larangan ini khusus untuk slide.
    showstringspaces=false,
    tabsize=4,
    columns=fullflexible,
    keepspaces=true,
    upquote=true,
    xleftmargin=12pt,
    aboveskip=0.4em,
    belowskip=0.4em
}

% ---------------------------------------------------------
% 3. Gaya keluaran program / sesi debugger untuk slide
% ---------------------------------------------------------
\lstdefinestyle{outslide}{
    basicstyle=\ttfamily\scriptsize,
    keywordstyle=\color{black},
    commentstyle=\color{obeGray},
    stringstyle=\color{black},
    numbers=none,
    backgroundcolor=\color{obeGrayLight},
    frame=single,
    rulecolor=\color{codeFrame},
    framesep=4pt,
    breaklines=true,
    showstringspaces=false,
    columns=fullflexible,
    keepspaces=true,
    upquote=true,
    xleftmargin=4pt,
    aboveskip=0.4em,
    belowskip=0.4em
}

\lstset{style=asmslide}

% ---------------------------------------------------------
% 4. Environment ringkas untuk dipakai di dalam frame
%
% Frame yang memuat kedua environment ini WAJIB memakai
% opsi [fragile], karena isinya verbatim.
% ---------------------------------------------------------
\lstnewenvironment{kodeasm}[1][]{%
    \lstset{style=asmslide, #1}%
}{}

\lstnewenvironment{keluaran}[1][]{%
    \lstset{style=outslide, #1}%
}{}



\renewcommand{\PertemuanTeks}{Pertemuan 9}

\begin{document}

% =========================================================
% 1. JUDUL
% =========================================================
\JudulDeck{Instruksi Transfer Data dan Aritmatika}{Pertemuan 9}

% =========================================================
% 2. CAPAIAN DAN TUJUAN
% =========================================================
\begin{frame}{Capaian dan Tujuan Pembelajaran}
\begin{block}{Sub-CPMK 3.1}
\small Menulis program assembly lengkap dengan instruksi transfer data
(\asm{MOV}, \asm{XCHG}, \asm{LEA}, \asm{PUSH}/\asm{POP}) dan instruksi
aritmatika (\asm{ADD}, \asm{SUB}, \asm{INC}, \asm{DEC}, \asm{MUL},
\asm{IMUL}, \asm{DIV}, \asm{IDIV}).
\end{block}

\vspace{0.1em}
{\small\textbf{CPMK yang didukung:} \CPMK{03}\quad\textbullet\quad
\textbf{CPL yang didukung:} \CPL{03}\par}

\vspace{0.05em}
\begin{exampleblock}{Setelah pertemuan ini, mahasiswa mampu:}
\footnotesize
\begin{enumerate}
    \item menerapkan aturan legalitas \asm{MOV} dan mengatasi pemindahan
          yang ditolak dengan register perantara;
    \item membedakan \asm{MOV} dan \asm{LEA}: nilai data atau alamat data;
    \item memakai \asm{PUSH} dan \asm{POP} untuk menyimpan dan mengambil
          kembali data pada \eng{stack};
    \item menjelaskan pengaruh aritmatika terhadap \flag{CF}, \flag{ZF},
          \flag{SF}, \flag{OF}, \flag{AF}, dan \flag{PF};
    \item menuliskan perkalian dan pembagian dengan operand implisit yang
          benar, menangani \eng{divide overflow}, dan menyusun aritmatika
          bertingkat dengan \asm{ADC} dan \asm{SBB}.
\end{enumerate}
\end{exampleblock}
\end{frame}

% =========================================================
% 3. PETA MATERI
% =========================================================
\begin{frame}{Peta Materi}
\small
\begin{enumerate}
    \item \textbf{Instruksi Transfer Data}\\
          {\footnotesize \asm{MOV} dan tiga aturan legalitasnya; register
          perantara; \asm{XCHG}; perbedaan \asm{MOV} dan \asm{LEA};
          \asm{PUSH}/\asm{POP}; \asm{LDS}/\asm{LES}}

    \vspace{0.25em}
    \item \textbf{Instruksi Aritmatika Dasar dan Pengaruhnya pada
          \eng{Flag}}\\
          {\footnotesize \asm{ADD} dan \asm{ADC}; \asm{SUB}, \asm{SBB},
          \asm{INC}, \asm{DEC}, \asm{NEG}; makna \eng{flag}; operasi
          bertanda dan tak bertanda; \asm{CMP}}

    \vspace{0.25em}
    \item \textbf{Perkalian, Pembagian, dan Aritmatika Bertingkat}\\
          {\footnotesize operand implisit; \eng{divide overflow};
          \asm{IMUL}/\asm{IDIV} serta \asm{CBW}/\asm{CWD}; \asm{ADC} dan
          \asm{SBB} untuk lebih dari 16 bit; konversi ke digit desimal}
\end{enumerate}

\vspace{0.3em}
\begin{block}{Praktik pada pertemuan ini}
\small Menelusuri \file{aritmatika.asm} dan \file{kali-bagi.asm}, mencatat
nilai \eng{flag}, lalu membangun kalkulator empat operasi sebagai Tugas~4.
\end{block}
\end{frame}

% =========================================================
% BAGIAN 1
% =========================================================
\PemisahBagian{Instruksi Transfer Data}

% ---------------------------------------------------------
\begin{frame}{Instruksi \asm{MOV}: Menyalin, Bukan Memindahkan}
\small \asm{MOV tujuan, sumber} menyalin data ke tujuan; isi sumber
\emph{tidak} terhapus. Setelah \asm{MOV AX, BX}, kedua register berisi
nilai yang sama.

\vspace{0.2em}
{\color{obeBlue}\bfseries Tiga aturan legalitas operand}
\vspace{0.1em}
\begin{enumerate}\footnotesize
    \item \textbf{Ukuran kedua operand harus sama.} \asm{MOV AX, BL}
          ditolak karena \reg{AX} 16 bit sedangkan \reg{BL} 8 bit.
    \item \textbf{Tidak boleh memori ke memori.} \asm{MOV nilai2, nilai1}
          ditolak; 8086 tidak memiliki jalur antarlokasi memori.
    \item \textbf{Tidak boleh nilai segera ke register segmen.}
          \asm{MOV DS, 1000h} ditolak; harus melalui register umum.
\end{enumerate}

\vspace{0.15em}
\begin{alertblock}{Jebakan yang sering terjadi}
\footnotesize Menambal \asm{MOV nilai2, nilai1} dengan \code{WORD PTR}
tidak menyelesaikan masalah --- penyebabnya jenis operand, bukan ukurannya.
\end{alertblock}
\Sumber{Bab X, Instruksi MOV dan Aturan Legalitasnya}
\end{frame}

% ---------------------------------------------------------
\begin{frame}[fragile]{Register Perantara: Cara Memperbaiki yang Ditolak}
\small Setiap kombinasi yang ditolak diselesaikan dengan menyisipkan
register perantara, bukan dengan menambah \code{PTR}.

\vspace{0.2em}
\begin{kodeasm}[numbers=none]
    ; memori ke memori : dua langkah melalui AX
    MOV AX, nilai1          ; muat isi memori ke AX
    MOV nilai2, AX          ; simpan dari AX ke memori tujuan

    ; nilai segera ke DS : dua langkah melalui AX
    MOV AX, 1000h           ; muat nilai segera ke AX
    MOV DS, AX              ; salin dari AX ke DS
\end{kodeasm}

\vspace{0.15em}
{\footnotesize Pola dua langkah yang sama juga dipakai untuk pemindahan
segmen ke segmen: \asm{MOV AX, ES} lalu \asm{MOV DS, AX}.\par}
\Sumber{Bab X, Mengatasi Pemindahan yang Tidak Diizinkan}
\end{frame}

% ---------------------------------------------------------
\begin{frame}{Tabel Legalitas Kombinasi Operand \asm{MOV}}
\footnotesize
\begin{center}
\begin{tabular}{@{}p{3.9cm}p{2.9cm}p{0.9cm}p{3.5cm}@{}}
    \toprule
    \textbf{Bentuk pemindahan} & \textbf{Contoh} & \textbf{Sah} &
    \textbf{Keterangan} \\
    \midrule
    Register ke register & \asm{MOV AX, BX} & Ya &
    Tanpa akses memori \\
    \addlinespace
    Nilai segera ke register & \asm{MOV AX, 5} & Ya &
    Nilai segera muat 16 bit \\
    \addlinespace
    Memori ke register & \asm{MOV AX, nilai1} & Ya &
    Satu operand di memori \\
    \addlinespace
    Register ke memori & \asm{MOV nilai1, AX} & Ya &
    Menyimpan hasil hitungan \\
    \addlinespace
    Register ke segmen & \asm{MOV DS, AX} & Ya &
    Cara baku inisialisasi \reg{DS} \\
    \addlinespace
    Nilai segera ke segmen & \asm{MOV DS, 1000h} & Tidak &
    Harus lewat register umum \\
    \addlinespace
    Memori ke memori & \asm{MOV nilai2, nilai1} & Tidak &
    Pakai register perantara \\
    \bottomrule
\end{tabular}
\end{center}
\Sumber{Bab X, Tabel Legalitas Kombinasi Operand pada Instruksi MOV}
\end{frame}

% ---------------------------------------------------------
\begin{frame}[fragile]{Instruksi \asm{XCHG}: Menukar Dua Operand}
\small \asm{XCHG} menukar isi dua operand sehingga keduanya berubah
sekaligus, tanpa memerlukan register ketiga.

\vspace{0.2em}
\begin{kodeasm}[numbers=none]
    MOV AX, 1234h       ; AX = 1234h
    MOV BX, 5678h       ; BX = 5678h
    XCHG AX, BX         ; AX = 5678h, BX = 1234h
\end{kodeasm}

\vspace{0.15em}
\begin{block}{Aturan \asm{XCHG} lebih ketat daripada \asm{MOV}}
\small Minimal satu operand wajib berada di register, kedua operand harus
berukuran sama, dan tidak ada operand yang boleh berupa nilai segera.
\end{block}
\Sumber{Bab X, Instruksi XCHG}
\end{frame}

% ---------------------------------------------------------
\begin{frame}[fragile]{Perbedaan \asm{MOV} dan \asm{LEA}}
\small \asm{LEA} memuat \emph{alamat} operand, sedangkan \asm{MOV} memuat
\emph{isi} operand.

\vspace{0.15em}
\begin{columns}[t]
\begin{column}{0.5\textwidth}
\begin{kodeasm}[numbers=none]
pesan DB 'Halo$'

    LEA BX, pesan     ; BX = alamat pesan
    MOV BX, OFFSET pesan
                      ; hasilnya sama
    MOV BX, pesan     ; BX = isi 'Ha'
\end{kodeasm}
\end{column}
\begin{column}{0.47\textwidth}
\begin{tikzpicture}[
    kotak/.style={draw=obeBlue, line width=0.8pt, fill=obeBlueLight,
        rounded corners=3pt, minimum width=2.6cm, minimum height=0.62cm,
        align=center, font=\scriptsize\ttfamily, inner sep=2pt},
    hasil/.style={draw=obeTeal, line width=0.9pt, fill=obeTealLight,
        rounded corners=3pt, minimum width=3.1cm, minimum height=0.62cm,
        align=center, font=\scriptsize, inner sep=2pt},
    panah/.style={-{Stealth[length=2mm]}, line width=0.9pt, obeTeal}
]
    \node[kotak] (a) at (0, 0.7) {LEA BX, pesan};
    \node[hasil] (b) at (0, -0.2) {BX = alamat \file{pesan}};
    \draw[panah] (a) -- (b);

    \node[kotak] (c) at (0, -1.1) {MOV BX, pesan};
    \node[hasil] (d) at (0, -2.0) {BX = isi \code{'Ha'}};
    \draw[panah] (c) -- (d);
\end{tikzpicture}
\end{column}
\end{columns}

\vspace{0.1em}
{\footnotesize\color{obeGray}Ketiganya sah secara sintaksis, sehingga
assembler tidak melaporkan kesalahan --- programnya saja yang salah.\par}
\end{frame}

% ---------------------------------------------------------
\begin{frame}{Instruksi \asm{PUSH} dan \asm{POP} pada \eng{Stack}}
\begin{columns}[t]
\begin{column}{0.53\textwidth}
\begin{center}
\begin{tikzpicture}[
    sel/.style={draw=obeBlue, line width=0.7pt, minimum width=1.7cm,
        minimum height=0.55cm, align=center, inner sep=1pt,
        font=\scriptsize},
    isi/.style={sel, fill=obeBlueLight},
    kosong/.style={sel, fill=obeGrayLight},
    sp/.style={-{Stealth[length=2mm]}, line width=1pt, obeMaroon},
    panah/.style={-{Stealth[length=2mm]}, line width=0.9pt, obeTeal},
    ket/.style={font=\bfseries\footnotesize, text=obeBlue}
]
    \node[ket] at (0, 2.35) {Sebelum};
    \node[kosong] at (0, 0) {kosong};
    \node[isi] at (0, 0.62) {5678h};
    \draw[sp] (0.95, 0.62) -- (0.87, 0.62);
    \node[font=\tiny, text=obeMaroon, anchor=west] at (1.0, 0.62) {SP};

    \node[ket] at (3.8, 2.35) {Sesudah};
    \node[sel, fill=obeAmberLight, draw=obeAmber] (ax) at (3.8, 1.55)
        {AX = 1234h};
    \node[isi] (atas) at (3.8, 0.62) {1234h};
    \node[isi] at (3.8, 0) {5678h};
    \draw[panah] (ax.south) -- (atas.north);
    \draw[sp] (4.75, 0.62) -- (4.67, 0.62);
    \node[font=\tiny, text=obeMaroon, anchor=west] at (4.8, 0.62) {SP};
\end{tikzpicture}

{\tiny\color{obeGray}\reg{SP} berkurang 2 setiap kali \asm{PUSH}.\par}
\end{center}
\end{column}
\begin{column}{0.45\textwidth}
{\footnotesize\color{obeBlue}\bfseries Aturan pemakaian}
\begin{itemize}\footnotesize
    \item Bekerja pada satuan word; \asm{PUSH AL} tidak tersedia.
    \item Nilai segera tidak diizinkan sebagai operand.
    \item \reg{CS} boleh di-\asm{PUSH}, tidak boleh menjadi tujuan
          \asm{POP}.
    \item Setiap \asm{PUSH} harus berpasangan dengan \asm{POP} yang
          seimbang agar \reg{SP} kembali ke posisi semula.
\end{itemize}
\end{column}
\end{columns}
\Sumber{Bab X, Instruksi PUSH dan POP}
\end{frame}

% ---------------------------------------------------------
\begin{frame}{Instruksi \asm{LDS}, \asm{LES}, dan Ringkasan Transfer Data}
{\footnotesize \asm{LDS} dan \asm{LES} memuat \eng{pointer} jauh 32 bit
sekaligus: offset ke register dan nomor segmen ke \reg{DS} atau
\reg{ES}.\par}

\vspace{0.2em}
{\scriptsize
\begin{center}
\begin{tabular}{@{}p{1.4cm}p{3.1cm}p{4.3cm}p{3.2cm}@{}}
    \toprule
    \textbf{Instruksi} & \textbf{Bentuk umum} & \textbf{Kegunaan} &
    \textbf{Contoh} \\
    \midrule
    \asm{MOV} & \asm{MOV tujuan, sumber} & Menyalin data; sumber tetap &
    \asm{MOV AX, BX} \\
    \addlinespace
    \asm{XCHG} & \asm{XCHG op1, op2} & Menukar isi dua operand &
    \asm{XCHG AX, BX} \\
    \addlinespace
    \asm{LEA} & \asm{LEA reg16, mem} & Memuat alamat efektif &
    \asm{LEA SI, pesan} \\
    \addlinespace
    \asm{LDS} & \asm{LDS reg16, mem32} & Offset ke reg, segmen ke \reg{DS} &
    \asm{LDS SI, ptrJauh} \\
    \addlinespace
    \asm{LES} & \asm{LES reg16, mem32} & Offset ke reg, segmen ke \reg{ES} &
    \asm{LES DI, ptrJauh} \\
    \addlinespace
    \asm{PUSH} & \asm{PUSH sumber} & Menaruh word ke \eng{stack} &
    \asm{PUSH AX} \\
    \addlinespace
    \asm{POP} & \asm{POP tujuan} & Mengambil word dari \eng{stack} &
    \asm{POP BX} \\
    \bottomrule
\end{tabular}
\end{center}
}
\Sumber{Bab X, Instruksi LDS dan LES serta Tabel Ringkasan Transfer Data}
\end{frame}

% ---------------------------------------------------------
\begin{frame}{Rangkuman Bagian 1}
\small
\begin{itemize}
    \item \asm{MOV} menyalin data tanpa menghapus sumber. Legalitasnya
          dibatasi tiga aturan: ukuran operand harus sama, tidak boleh
          memori ke memori, dan tidak boleh nilai segera ke register
          segmen.

    \vspace{0.22em}
    \item Pemindahan yang ditolak diselesaikan dengan register perantara
          dua langkah, bukan dengan menambahkan \code{WORD PTR}.

    \vspace{0.22em}
    \item \asm{XCHG} menukar isi dua operand dan mensyaratkan minimal satu
          operand berada di register; nilai segera tidak diizinkan.

    \vspace{0.22em}
    \item \asm{MOV} memindahkan \emph{isi} memori, sedangkan \asm{LEA}
          memindahkan \emph{alamat} memori --- lebih ringkas dan tidak
          menyentuh memori bila alamat dihitung saat program berjalan.

    \vspace{0.22em}
    \item \asm{PUSH} dan \asm{POP} bekerja pada satuan word, tidak
          menerima nilai segera, dan harus berpasangan seimbang agar
          \reg{SP} kembali ke posisi semula.
\end{itemize}
\end{frame}

% =========================================================
% BAGIAN 2
% =========================================================
\PemisahBagian{Instruksi Aritmatika Dasar dan Pengaruhnya pada \eng{Flag}}

% ---------------------------------------------------------
\begin{frame}[fragile]{Penjumlahan: \asm{ADD} dan \asm{ADC}}
\small \asm{ADD} menjumlahkan sumber ke tujuan dan menyimpan hasilnya pada
tujuan. \asm{ADC} bekerja sama, tetapi ikut menambahkan \flag{CF} dari
operasi sebelumnya.

\vspace{0.2em}
\begin{kodeasm}[numbers=none]
    MOV AX, 25              ; AX = 25
    ADD AX, 17              ; AX = 42
    ADD nilai1, AX          ; nilai1 = nilai1 + AX (di memori)
    ADC AX, BX              ; AX = AX + BX + CF
\end{kodeasm}

\vspace{0.15em}
\begin{block}{Kapan perbedaannya terasa}
\small Pada operasi satu langkah keduanya sering dianggap sama karena
\flag{CF} biasanya nol. Bedanya baru tampak pada aritmatika bertingkat,
yaitu saat \eng{carry} harus diteruskan ke word berikutnya.
\end{block}
\Sumber{Bab X, Penjumlahan ADD dan ADC}
\end{frame}

% ---------------------------------------------------------
\begin{frame}[fragile]{Pengurangan dan Penyangkalan}
\small \asm{SUB} mengurangi tujuan dengan sumber, \asm{SBB} ikut
mengurangkan \flag{CF}, \asm{INC} dan \asm{DEC} menambah atau mengurangi
satu, dan \asm{NEG} mengubah tanda bilangan.

\vspace{0.2em}
\begin{kodeasm}[numbers=none]
    MOV AX, 100             ; AX = 100
    SUB AX, 45              ; AX = 55
    INC AX                  ; AX = 56
    DEC AX                  ; AX = 55
    NEG AX                  ; AX = -55 (0FFC9h)
    SBB AX, BX              ; AX = AX - BX - CF
\end{kodeasm}

\vspace{0.15em}
\begin{alertblock}{Dua kesalahan yang sering terjadi}
\small \asm{NEG} tidak selalu menyetel \flag{CF}: bila \reg{AX} nol,
hasilnya nol dan \flag{CF} bernilai 0. \asm{INC} sengaja tidak mengubah
\flag{CF}, sehingga tidak boleh dipakai pada aritmatika bertingkat.
\end{alertblock}
\Sumber{Bab X, Pengurangan dan Penyangkalan}
\end{frame}

% ---------------------------------------------------------
\begin{frame}[fragile]{Perbandingan dengan \asm{CMP}}
\small \asm{CMP} bekerja identik dengan \asm{SUB}, tetapi \emph{tidak
menyimpan} hasilnya. Yang ditinggalkan hanyalah jejak pada \eng{flag},
sehingga \asm{CMP} selalu diikuti lompatan bersyarat.

\vspace{0.2em}
\begin{columns}[t]
\begin{column}{0.52\textwidth}
\begin{kodeasm}[numbers=none]
    CMP AX, BX         ; AX - BX
    JE  SAMA           ; ZF = 1
    JA  LEBIH_BESAR    ; CF = 0, ZF = 0
    JG  LEBIH_TTD      ; SF = OF, ZF = 0
\end{kodeasm}
\end{column}
\begin{column}{0.45\textwidth}
\begin{exampleblock}{Tiga pola yang paling sering}
\footnotesize
\asm{JE} untuk kesamaan; \asm{JA} untuk tak bertanda; \asm{JG} untuk
bertanda. Pilihan lompatan harus sejalan dengan tafsiran operandnya.
\end{exampleblock}
\end{column}
\end{columns}

\vspace{0.15em}
{\footnotesize\color{obeGray}Karena hasilnya dibuang, \asm{CMP} tidak
mengubah nilai register apa pun --- hanya \eng{flag}.\par}
\Sumber{Bab X, Operasi Perbandingan dengan CMP}
\end{frame}

% ---------------------------------------------------------
\begin{frame}{Makna \eng{Flag} Hasil Aritmatika}
\scriptsize
\begin{center}
\begin{tabular}{@{}p{0.9cm}p{6.4cm}p{5.0cm}@{}}
    \toprule
    \textbf{Flag} & \textbf{Arti} & \textbf{Tanda pada hasil} \\
    \midrule
    \flag{CF} & \eng{Carry}: ada nilai yang keluar dari bit paling kiri;
    pada pengurangan berarti pinjaman &
    \flag{CF} = 1 menandakan hasil tak bertanda melampaui kapasitas \\
    \addlinespace
    \flag{ZF} & \eng{Zero}: hasil operasi bernilai nol &
    \flag{ZF} = 1 menandakan kedua operand sama besar \\
    \addlinespace
    \flag{SF} & \eng{Sign}: bit paling kiri hasil bernilai 1 &
    \flag{SF} = 1 menandakan hasil bertanda negatif \\
    \addlinespace
    \flag{OF} & \eng{Overflow}: hasil bertanda melampaui batas yang dapat
    ditampung & \flag{OF} = 1 menandakan luapan bilangan bertanda \\
    \addlinespace
    \flag{AF} & \eng{Auxiliary carry} dari bit ke-3 menuju bit ke-4 &
    Dipakai instruksi penyesuaian BCD seperti \asm{DAA} \\
    \addlinespace
    \flag{PF} & \eng{Parity}: cacah bit 1 pada byte paling rendah genap &
    Dipakai pada pemeriksaan kesalahan transmisi \\
    \bottomrule
\end{tabular}
\end{center}
\Sumber{Bab X, Makna Flag Hasil Aritmatika}
\end{frame}

% ---------------------------------------------------------
\begin{frame}{Ringkasan Instruksi Aritmatika dan \eng{Flag}-nya}
\scriptsize
\begin{center}
\begin{tabular}{@{}p{1.4cm}p{2.6cm}p{4.1cm}p{3.8cm}@{}}
    \toprule
    \textbf{Instruksi} & \textbf{Bentuk operand} & \textbf{Hasil} &
    \textbf{Flag terpengaruh} \\
    \midrule
    \asm{ADD} & \asm{ADD tujuan, sumber} & tujuan $\leftarrow$ tujuan $+$
    sumber & \flag{OF} \flag{SF} \flag{ZF} \flag{AF} \flag{PF} \flag{CF} \\
    \addlinespace
    \asm{ADC} & \asm{ADC tujuan, sumber} & tujuan $\leftarrow$ tujuan $+$
    sumber $+$ \flag{CF} & seluruh \eng{flag} aritmatika \\
    \addlinespace
    \asm{SUB} & \asm{SUB tujuan, sumber} & tujuan $\leftarrow$ tujuan $-$
    sumber & seluruh \eng{flag} aritmatika \\
    \addlinespace
    \asm{SBB} & \asm{SBB tujuan, sumber} & tujuan $\leftarrow$ tujuan $-$
    sumber $-$ \flag{CF} & seluruh \eng{flag} aritmatika \\
    \addlinespace
    \asm{INC} & \asm{INC tujuan} & tujuan $\leftarrow$ tujuan $+$ 1 &
    \flag{OF} \flag{SF} \flag{ZF} \flag{AF} \flag{PF}; \flag{CF} tetap \\
    \addlinespace
    \asm{DEC} & \asm{DEC tujuan} & tujuan $\leftarrow$ tujuan $-$ 1 &
    \flag{OF} \flag{SF} \flag{ZF} \flag{AF} \flag{PF}; \flag{CF} tetap \\
    \addlinespace
    \asm{NEG} & \asm{NEG tujuan} & tujuan $\leftarrow$ 0 $-$ tujuan &
    seluruh \eng{flag}; \flag{CF} = 1 bila hasil tidak nol \\
    \addlinespace
    \asm{CMP} & \asm{CMP tujuan, sumber} & hasil dibuang, hanya \eng{flag}
    yang diubah & sama dengan \asm{SUB} \\
    \bottomrule
\end{tabular}
\end{center}
\Sumber{Bab X, Ringkasan Instruksi Aritmatika 8086}
\end{frame}

% ---------------------------------------------------------
\begin{frame}[fragile]{Operasi Bertanda dan Tak Bertanda}
\small Satu deretan bit yang sama dapat ditafsirkan sebagai bilangan tak
bertanda atau bertanda dalam komplemen dua. Byte \hex{FF} berarti 255 bila
tak bertanda, tetapi $-1$ bila bertanda.

\vspace{0.15em}
\begin{columns}[t]
\begin{column}{0.5\textwidth}
\begin{kodeasm}[numbers=none]
    ; tak bertanda : 200 + 100
    MOV AL, 200     ; AL = C8h
    ADD AL, 100     ; AL = 2Ch, CF = 1
                    ; -> salah

    ; bertanda : -56 + 100
    MOV AL, 200     ; AL = C8h = -56
    ADD AL, 100     ; AL = 2Ch, OF = 0
                    ; -> benar
\end{kodeasm}
\end{column}
\begin{column}{0.47\textwidth}
\begin{block}{Aturan praktis}
\footnotesize
\flag{CF} untuk hasil tak bertanda dan \flag{OF} untuk hasil bertanda.
Karena itu, pakai lompatan tak bertanda (\asm{JA}, \asm{JB}) setelah
operasi tak bertanda, dan lompatan bertanda (\asm{JG}, \asm{JL}) setelah
operasi bertanda.
\end{block}
{\footnotesize\color{obeGray}Perbedaan yang sama mendasari pemilihan
\asm{MUL} atau \asm{IMUL}.\par}
\end{column}
\end{columns}
\end{frame}

% ---------------------------------------------------------
\begin{frame}{Rangkuman Bagian 2}
\small
\begin{itemize}
    \item Instruksi aritmatika dasar: \asm{ADD}, \asm{ADC}, \asm{SUB},
          \asm{SBB}, \asm{INC}, \asm{DEC}, \asm{NEG}, dan \asm{CMP}.

    \vspace{0.22em}
    \item Setiap instruksi meninggalkan jejak pada \eng{flag}, kecuali
          \flag{CF} yang tidak disentuh \asm{INC} dan \asm{DEC} --- itulah
          sebabnya keduanya tidak dipakai pada aritmatika bertingkat.

    \vspace{0.22em}
    \item \flag{CF} menandai luapan tak bertanda, \flag{OF} menandai
          luapan bertanda. Pilihan instruksi dan lompatan bersyarat harus
          mengikuti tafsiran yang dipilih.

    \vspace{0.22em}
    \item \asm{NEG} menghitung $0 - AX$, sehingga \flag{CF} bernilai 0
          ketika hasilnya nol.

    \vspace{0.22em}
    \item \asm{CMP} tidak menyimpan hasil, hanya menyetel \eng{flag} sama
          seperti \asm{SUB}, lalu diikuti \asm{JE}, \asm{JA}, atau
          \asm{JG}.
\end{itemize}
\end{frame}

% =========================================================
% BAGIAN 3
% =========================================================
\PemisahBagian{Perkalian, Pembagian, dan Aritmatika Bertingkat}

% ---------------------------------------------------------
\begin{frame}{Aturan Operand Implisit}
\small \asm{MUL}, \asm{IMUL}, \asm{DIV}, dan \asm{IDIV} hanya menuliskan
satu operand. Operand satunya sudah ditentukan prosesor dan tidak ditulis
di dalam instruksi.

\vspace{0.15em}
\begin{center}
\begin{tikzpicture}[
    kotak/.style={draw=obeBlue, line width=0.8pt, fill=obeBlueLight,
        minimum width=1.3cm, minimum height=0.55cm, align=center,
        inner sep=2pt, font=\scriptsize},
    hasil/.style={draw=obeAmber, line width=0.8pt, fill=obeAmberLight,
        minimum width=1.7cm, minimum height=0.55cm, align=center,
        inner sep=2pt, font=\scriptsize},
    op/.style={draw=obeTeal, line width=1pt, fill=obeTealLight,
        rounded corners=2pt, minimum width=1.1cm, minimum height=0.55cm,
        align=center, font=\scriptsize\bfseries, inner sep=2pt},
    panah/.style={-{Stealth[length=1.8mm]}, line width=0.8pt, obeGray},
    baris/.style={font=\scriptsize\bfseries, text=obeBlue, align=right,
        anchor=east},
    ket/.style={font=\tiny, text=obeGray, align=left, anchor=west}
]
    \node[baris] at (-0.15, 2.7) {\asm{MUL} r/m8};
    \node[kotak] (a1) at (1.2, 3.0) {AL};
    \node[kotak] (a2) at (1.2, 2.4) {r/m8};
    \node[op]    (a3) at (3.3, 2.7) {MUL};
    \node[hasil] (a4) at (5.7, 2.7) {AX};
    \draw[panah] (a1.east) -- (a3.west);
    \draw[panah] (a2.east) -- (a3.west);
    \draw[panah] (a3.east) -- (a4.west);

    \node[baris] at (-0.15, 1.55) {\asm{MUL} r/m16};
    \node[kotak] (b1) at (1.2, 1.85) {AX};
    \node[kotak] (b2) at (1.2, 1.25) {r/m16};
    \node[op]    (b3) at (3.3, 1.55) {MUL};
    \node[hasil] (b4) at (5.7, 1.55) {DX:AX};
    \draw[panah] (b1.east) -- (b3.west);
    \draw[panah] (b2.east) -- (b3.west);
    \draw[panah] (b3.east) -- (b4.west);

    \node[baris] at (-0.15, 0.4) {\asm{DIV} r/m8};
    \node[kotak] (c1) at (1.2, 0.7) {AX};
    \node[kotak] (c2) at (1.2, 0.1) {r/m8};
    \node[op]    (c3) at (3.3, 0.4) {DIV};
    \node[hasil] (c4) at (5.7, 0.7) {AL};
    \node[hasil] (c5) at (5.7, 0.1) {AH};
    \draw[panah] (c1.east) -- (c3.west);
    \draw[panah] (c2.east) -- (c3.west);
    \draw[panah] (c3.east) -- (c4.west);
    \draw[panah] (c3.east) -- (c5.west);
    \node[ket] at (6.68, 0.7) {hasil bagi};
    \node[ket] at (6.68, 0.1) {sisa};

    \node[baris] at (-0.15, -0.75) {\asm{DIV} r/m16};
    \node[kotak] (d1) at (1.2, -0.45) {DX:AX};
    \node[kotak] (d2) at (1.2, -1.05) {r/m16};
    \node[op]    (d3) at (3.3, -0.75) {DIV};
    \node[hasil] (d4) at (5.7, -0.45) {AX};
    \node[hasil] (d5) at (5.7, -1.05) {DX};
    \draw[panah] (d1.east) -- (d3.west);
    \draw[panah] (d2.east) -- (d3.west);
    \draw[panah] (d3.east) -- (d4.west);
    \draw[panah] (d3.east) -- (d5.west);
    \node[ket] at (6.68, -0.45) {hasil bagi};
    \node[ket] at (6.68, -1.05) {sisa};
\end{tikzpicture}
\end{center}
\Sumber{Bab X, Aturan Operand Implisit}
\end{frame}

% ---------------------------------------------------------
\begin{frame}{Operand Implisit dan Ukuran Hasil}
\scriptsize
\begin{center}
\begin{tabular}{@{}p{1.4cm}p{2.0cm}p{3.7cm}p{4.5cm}@{}}
    \toprule
    \textbf{Instruksi} & \textbf{Operand} & \textbf{Yang dihitung} &
    \textbf{Letak hasil} \\
    \midrule
    \asm{MUL} & r/m8 & \reg{AL} $\times$ operand &
    \reg{AX}; \flag{CF} = \flag{OF} = 1 bila \reg{AH} tidak nol \\
    \addlinespace
    \asm{MUL} & r/m16 & \reg{AX} $\times$ operand &
    \reg{DX}:\reg{AX}; \flag{CF} = \flag{OF} = 1 bila \reg{DX} tidak nol \\
    \addlinespace
    \asm{IMUL} & r/m8 & \reg{AL} $\times$ operand (bertanda) &
    \reg{AX}; \flag{OF} = 1 bila hasil tidak muat di \reg{AL} \\
    \addlinespace
    \asm{IMUL} & r/m16 & \reg{AX} $\times$ operand (bertanda) &
    \reg{DX}:\reg{AX}; \flag{OF} = 1 bila tidak muat di \reg{AX} \\
    \addlinespace
    \asm{DIV} & r/m8 & \reg{AX} $\div$ operand &
    hasil bagi di \reg{AL}, sisa di \reg{AH} \\
    \addlinespace
    \asm{DIV} & r/m16 & \reg{DX}:\reg{AX} $\div$ operand &
    hasil bagi di \reg{AX}, sisa di \reg{DX} \\
    \bottomrule
\end{tabular}
\end{center}

\vspace{0.2em}
{\footnotesize Hasil perkalian selalu berukuran dua kali operand, sehingga
\asm{MUL} tidak pernah kehilangan data; yang ada hanyalah penanda
\flag{CF} dan \flag{OF}.\par}
\Sumber{Bab X, Operand Implisit dan Ukuran Hasil}
\end{frame}

% ---------------------------------------------------------
\begin{frame}[fragile]{Kesalahan \eng{Divide Overflow}}
\small Pembagian berperilaku sebaliknya: hasil bagi harus muat pada
register tujuan yang ukurannya hanya separuh operand. Bila tidak muat,
prosesor membangkitkan kesalahan \eng{divide overflow} dan program DOS
umumnya berhenti mendadak.

\vspace{0.2em}
\begin{columns}[t]
\begin{column}{0.55\textwidth}
\begin{kodeasm}[numbers=none]
    MOV AX, 1000    ; AX = 03E8h
    MOV BL, 7       ; BL = pembagi
    CMP AH, BL      ; luapan bila AH >= BL
    JAE LUAPAN      ; periksa sebelum DIV
    DIV BL          ; AL = 142, AH = 6

LUAPAN:
    ...             ; laporkan, jangan DIV
\end{kodeasm}
\end{column}
\begin{column}{0.42\textwidth}
\begin{alertblock}{Syarat aman}
\footnotesize
\asm{DIV BL} aman hanya bila \reg{AH} lebih kecil daripada \reg{BL}.
Untuk pembagian 16 bit, syaratnya \reg{DX} lebih kecil daripada operand
pembagi.
\end{alertblock}
{\footnotesize\color{obeGray}Luapan pembagian bukan kesalahan sintaks,
sehingga tidak dilaporkan assembler.\par}
\end{column}
\end{columns}
\Sumber{Bab X, Kesalahan Divide Overflow}
\end{frame}

% ---------------------------------------------------------
\begin{frame}[fragile]{Perkalian dan Pembagian Bertanda}
\small \asm{IMUL} dan \asm{IDIV} menafsirkan operand sebagai bilangan
bertanda dalam komplemen dua. Sebelum \asm{IDIV}, bilangan yang dibagi
harus diperluas tandanya.

\vspace{0.2em}
\begin{columns}[t]
\begin{column}{0.5\textwidth}
\begin{kodeasm}[numbers=none]
    MOV AX, 1000    ; AX = 03E8h
    MOV BX, 7       ; BX = 0007h
    CWD             ; DX = 0000h
    IDIV BX         ; AX = 142, DX = 6
\end{kodeasm}
\end{column}
\begin{column}{0.47\textwidth}
\begin{block}{\asm{CBW} dan \asm{CWD}}
\footnotesize
\asm{CBW} memperluas tanda \reg{AL} ke \reg{AX}; \asm{CWD} memperluas
tanda \reg{AX} ke \reg{DX}:\reg{AX}.
\end{block}
\end{column}
\end{columns}

\vspace{0.15em}
{\footnotesize\color{obeGray}Kesalahan karena \asm{CWD} dihilangkan tidak
terlihat pada bilangan positif; pada bilangan negatif, \reg{DX} seharusnya
bernilai \hex{FFFF} agar tafsiran bertanda menjadi benar.\par}
\Sumber{Bab X, Perkalian dan Pembagian Bertanda}
\end{frame}

% ---------------------------------------------------------
\begin{frame}[fragile]{Aritmatika Bertingkat dengan \asm{ADC} dan \asm{SBB}}
\small Bilangan 32 bit dipecah menjadi word rendah (bit 0--15) dan word
tinggi (bit 16--31). Penjumlahan dimulai dari word rendah; word tinggi
memakai \asm{ADC} agar \eng{carry} ikut diperhitungkan.

\vspace{0.15em}
\begin{columns}[t]
\begin{column}{0.5\textwidth}
\begin{kodeasm}[numbers=none]
    MOV AX, rendahA     ; word rendah A
    ADD AX, rendahB     ; AX = 1356h, CF = 1
    MOV hasilL, AX

    MOV AX, tinggiA     ; word tinggi A
    ADC AX, tinggiB     ; AX = 8ACFh
    MOV hasilT, AX
\end{kodeasm}
\end{column}
\begin{column}{0.47\textwidth}
\begin{exampleblock}{Propagasi \eng{carry}}
\footnotesize
Word rendah: \hex{CDEF} $+$ \hex{4567} $=$ \hex{1356}, maka
\flag{CF} = 1.\\[0.3em]
Word tinggi: \hex{89AB} $+$ \hex{0123} $+$ 1 $=$ \hex{8ACF}.\\[0.3em]
Hasil akhir: \hex{8ACF1356}.
\end{exampleblock}

\vspace{0.25em}
{\footnotesize\color{obeGray}Bila instruksi kedua ditulis \asm{ADD},
hasilnya \hex{8ACE} dan satu nilai hilang. Pengurangan 32 bit memakai pola
\asm{SUB} dan \asm{SBB}.\par}
\end{column}
\end{columns}
\Sumber{Bab X, Aritmatika Bertingkat dengan ADC dan SBB}
\end{frame}

% ---------------------------------------------------------
\begin{frame}[fragile]{Mengubah Bilangan Biner Menjadi Digit Desimal}
\small Nilai di register selalu biner, sedangkan \code{INT 21h} fungsi
\hex{02} mencetak satu karakter pada satu waktu. Bilangan harus dipecah
menjadi digit dengan pembagian berulang sepuluh.

\vspace{0.15em}
\begin{columns}[t]
\begin{column}{0.52\textwidth}
\begin{kodeasm}[numbers=none]
    MOV BX, 10          ; pembagi desimal
BAGI_DIGIT:
    XOR DX, DX          ; siapkan DX:AX
    DIV BX              ; AX bagi, DX sisa
    PUSH DX             ; simpan digit
    INC CX
    CMP AX, 0
    JNE BAGI_DIGIT
CETAK_DIGIT:
    POP DX
    ADD DL, 30h         ; nilai jadi ASCII
    MOV AH, 02h
    INT 21h
    LOOP CETAK_DIGIT
\end{kodeasm}
\end{column}
\begin{column}{0.45\textwidth}
{\footnotesize\color{obeBlue}\bfseries Alur algoritma}
\begin{enumerate}\footnotesize
    \item Bagi nilai dengan 10; sisanya satu digit.
    \item Simpan sisa pada \eng{stack}, tambah cacah digit.
    \item Ulangi selama hasil bagi masih lebih besar dari nol.
    \item Ambil kembali dari \eng{stack} agar urutan cetak terbalik
          menjadi benar.
    \item Tambahkan \hex{30} agar nilai $0$--$9$ menjadi kode ASCII
          \code{'0'}--\code{'9'}.
\end{enumerate}
\end{column}
\end{columns}
\Sumber{Bab X, Mengubah Bilangan Biner Menjadi Digit Desimal}
\end{frame}

% ---------------------------------------------------------
\begin{frame}{Studi Kasus: Kalkulator Empat Operasi Sederhana}
\small Kerangka program kalkulator mencakup tiga hal sekaligus: masukan
pengguna, pemilihan operasi, dan pencetakan hasil.

\vspace{0.3em}
\begin{enumerate}
    \item Tampilkan menu: \code{1} tambah, \code{2} kurang, \code{3} kali,
          \code{4} bagi.

    \vspace{0.2em}
    \item Baca pilihan, lalu arahkan alur program dengan \asm{CMP} dan
          lompatan bersyarat.

    \vspace{0.2em}
    \item Baca dua bilangan satu digit.

    \vspace{0.2em}
    \item Jalankan operasi yang dipilih: \asm{ADD}, \asm{SUB}, \asm{MUL},
          atau \asm{DIV}.

    \vspace{0.2em}
    \item Untuk pembagian, periksa pembagi nol dan hasil bagi yang tidak
          muat sebelum \asm{DIV} dijalankan.

    \vspace{0.2em}
    \item Cetak hasil dengan prosedur konversi desimal \code{CETAK\_AX}.
\end{enumerate}

\vspace{0.2em}
{\footnotesize\color{obeGray}Kerangka ini menjadi dasar Tugas~4 pada
bagian praktikum.\par}
\Sumber{Bab X, Studi Kasus Kalkulator Empat Operasi}
\end{frame}

% ---------------------------------------------------------
\begin{frame}{Rangkuman Bagian 3}
\small
\begin{itemize}
    \item \asm{MUL}, \asm{IMUL}, \asm{DIV}, dan \asm{IDIV} memakai operand
          implisit \reg{AL} atau \reg{AX}, dan menaruh hasil pada
          \reg{AX} atau \reg{DX}:\reg{AX}.

    \vspace{0.22em}
    \item Hasil perkalian selalu dua kali ukuran operand, sehingga
          \asm{MUL} tidak pernah kehilangan data; yang muncul hanya
          \flag{CF} dan \flag{OF}.

    \vspace{0.22em}
    \item \asm{DIV} menuntut hasil bagi muat pada register tujuan. Syarat
          amannya \reg{AH} $<$ \reg{BL}, diperiksa dengan
          \asm{CMP AH, BL} sebelum \asm{DIV BL} dijalankan.

    \vspace{0.22em}
    \item Sebelum \asm{IDIV}, tanda harus diperluas dengan \asm{CBW} atau
          \asm{CWD}; luapan pembagian tidak terdeteksi saat perakitan.

    \vspace{0.22em}
    \item Aritmatika bertingkat memakai \asm{ADC} dan \asm{SBB} agar
          \eng{carry} dan pinjaman diteruskan dari word rendah ke word
          tinggi; konversi desimal memakai pembagian berulang sepuluh.
\end{itemize}
\end{frame}

% =========================================================
% CONTOH PROGRAM
% =========================================================
\begin{frame}[fragile]{Contoh 1: Kalkulator Penjumlahan dan Pengurangan}
\begin{columns}[t]
\begin{column}{0.56\textwidth}
\begin{kodeasm}[numbers=none]
    MOV AH, 01h         ; baca satu digit
    INT 21h
    SUB AL, 30h         ; ASCII -> nilai 0-9
    XOR AH, AH          ; nolkan AH
    MOV nilai1, AX      ; simpan bilangan 1
    ...
    MOV AX, nilai1
    ADD AX, nilai2      ; hasil tambah
    MOV BX, AX          ; simpan agar AX bebas
    MOV AX, nilai1
    SUB AX, nilai2      ; hasil kurang
    JNS CETAK_KURANG
    NEG AX              ; jadikan nilai mutlak
CETAK_KURANG:
    CALL CETAK_AX       ; cetak AX desimal
\end{kodeasm}
\end{column}
\begin{column}{0.41\textwidth}
{\color{obeTeal}\bfseries Keluaran}\\[0.1em]
\begin{keluaran}
Masukkan angka pertama (0-9): 3
Masukkan angka kedua   (0-9): 8
Hasil tambah : 11
Hasil kurang : -5
\end{keluaran}

\vspace{0.3em}
{\footnotesize Hasil tambah disimpan di \reg{BX} karena \reg{AX} akan
dipakai prosedur pencetakan. Prosedur \code{CETAK\_AX} membagi AX dengan
10 berulang kali dan menyimpan sisanya pada \eng{stack}.\par}

\vspace{0.25em}
{\footnotesize\color{obeGray}\file{aritmatika.asm}\par}
\end{column}
\end{columns}
\end{frame}

% ---------------------------------------------------------
\begin{frame}[fragile]{Contoh 2: Perkalian, Pembagian, dan \eng{Divide Overflow}}
\begin{columns}[t]
\begin{column}{0.55\textwidth}
\begin{kodeasm}[numbers=none]
    MOV AL, 25      ; pengali 8 bit
    MOV BL, 4
    MUL BL          ; AX = AL x BL = 100

    MOV AL, 0FDh    ; 0FDh = -3
    MOV BL, 5
    IMUL BL         ; AX = -15 (0FFF1h)
    ...
    MOV AX, 1000    ; bilangan yang dibagi
    MOV BL, 7       ; pembagi
    CMP AH, BL      ; luapan bila AH >= BL
    JAE LUAPAN      ; periksa sebelum DIV
    DIV BL          ; AL = 142, AH = 6
\end{kodeasm}
\end{column}
\begin{column}{0.42\textwidth}
{\color{obeTeal}\bfseries Keluaran}\\[0.1em]
\begin{keluaran}
MUL  : 25 x 4   = 100
IMUL : -3 x 5   = -15
DIV  : 1000 / 7 = 142  sisa 6
DIV  : 1000 / 3 = divide
overflow: hasil bagi tidak
muat pada AL
\end{keluaran}

\vspace{0.25em}
{\footnotesize\color{obeGray}$1000 \div 3 = 333$ tidak muat pada
\reg{AL} yang hanya menampung $0$--$255$. Tanpa pemeriksaan, program
berhenti mendadak; kesalahan ini tidak terdeteksi saat
perakitan.\par}

\vspace{0.2em}
{\footnotesize\color{obeGray}\file{kali-bagi.asm}\par}
\end{column}
\end{columns}
\end{frame}

% ---------------------------------------------------------
\begin{frame}[fragile]{Contoh 3: Penjumlahan 32 Bit dengan \asm{ADC}}
\begin{columns}[t]
\begin{column}{0.56\textwidth}
\begin{kodeasm}[numbers=none]
    rendahA DW 0CDEFh   ; word rendah A
    tinggiA DW 89ABh    ; word tinggi A
    rendahB DW 4567h    ; word rendah B
    tinggiB DW 0123h    ; word tinggi B
    ...
    MOV AX, rendahA     ; AX = CDEFh
    ADD AX, rendahB     ; AX = 1356h, CF = 1
    MOV hasilL, AX      ; simpan word rendah

    MOV AX, tinggiA     ; AX = 89ABh
    ADC AX, tinggiB     ; AX = 8ACFh
    MOV hasilT, AX      ; simpan word tinggi

    MOV AH, 4Ch
    INT 21h
\end{kodeasm}
\end{column}
\begin{column}{0.41\textwidth}
\begin{exampleblock}{Yang diperlihatkan}
\footnotesize
\asm{ADD} pada tahap pertama menghasilkan word rendah sekaligus
meninggalkan \flag{CF} = 1. Nilai inilah yang ``dibawa'' ke tahap kedua
oleh \asm{ADC}.
\end{exampleblock}

\vspace{0.3em}
{\footnotesize Hasil akhirnya \hex{8ACF1356}. Pola yang sama berlaku
untuk pengurangan 32 bit dengan pasangan \asm{SUB} dan \asm{SBB}.\par}

\vspace{0.25em}
{\footnotesize\color{obeGray}Penerapan langsung dari tabel propagasi
\eng{carry}.\par}
\end{column}
\end{columns}
\end{frame}

% =========================================================
% PRAKTIKUM
% =========================================================
\begin{frame}{Praktikum: Membangun Kalkulator Sederhana}
\begin{block}{Judul praktikum}
\small Membangun kalkulator sederhana (masukan--proses--keluaran)
\quad\textbullet\quad \textbf{Sub-CPMK:} 3.1
\quad\textbullet\quad \textbf{Waktu:} 100 menit
\quad\textbullet\quad \textbf{Bentuk:} individu \quad\textbullet\quad
\textbf{Tugas 4} dikumpulkan
\end{block}

\vspace{0.2em}
\begin{columns}[t]
\begin{column}{0.48\textwidth}
{\color{obeBlue}\bfseries Alat dan bahan}
\begin{itemize}\small
    \item Komputer dengan Emu8086, GUI TASM, atau DOSBox
    \item Berkas \file{aritmatika.asm} dan \file{kali-bagi.asm}
    \item Lembar kerja Tugas~4 dan jurnal praktikum
\end{itemize}
\end{column}
\begin{column}{0.48\textwidth}
{\color{obeBlue}\bfseries Langkah kerja}
\begin{enumerate}\footnotesize
    \item Rakit dan jalankan kedua contoh; catat keluarannya.
    \item Amati \reg{AX}, \reg{BX}, \flag{CF}, \flag{ZF}, \flag{OF}, dan
          \reg{SP} dengan \eng{Single Step}.
    \item Hapus sementara \asm{CMP AH, BL}, lalu amati $1000 \div 3$.
    \item Salin kerangka menjadi \file{kalkulator.asm} dan tambahkan menu
          empat operasi.
\end{enumerate}
\end{column}
\end{columns}

\vspace{0.05em}
{\footnotesize\textbf{Dikumpulkan:} \file{kalkulator.asm}, tangkapan layar
keempat operasi dan pesan pembagi nol, tabel pengamatan \eng{flag}, serta
jurnal praktikum.\par}
\end{frame}

% =========================================================
% LATIHAN
% =========================================================
\begin{frame}{Latihan}
\small
\begin{enumerate}
    \item Jelaskan tiga aturan legalitas \asm{MOV}. Untuk setiap aturan,
          berikan satu contoh yang sah dan satu contoh yang ditolak
          assembler.

    \vspace{0.25em}
    \item Apakah perbedaan hasil antara \asm{LEA BX, pesan} dan
          \asm{MOV BX, pesan}? Kapan \asm{LEA} lebih tepat dipakai
          daripada \asm{MOV} dengan \asm{OFFSET}?

    \vspace{0.25em}
    \item Diberikan \asm{MOV AL, 250} lalu \asm{ADD AL, 15}. Berapa nilai
          akhir \reg{AL} dalam heksadesimal, dan berapa nilai \flag{CF},
          \flag{ZF}, \flag{SF}, serta \flag{OF}?

    \vspace{0.25em}
    \item Sebuah program menjalankan \asm{DIV BL} dengan \reg{AX} =
          \hex{0BB8} dan \reg{BL} = \hex{03}, lalu berhenti mendadak.
          Jelaskan penyebabnya dan tuliskan pemeriksaan yang perlu
          ditambahkan.

    \vspace{0.25em}
    \item Tulislah program yang menjumlahkan $A = \hex{0001FFFF}$ dan
          $B = \hex{00020001}$ dengan \asm{ADD} dan \asm{ADC}.
          Tunjukkan word rendah, word tinggi, dan \flag{CF} setiap tahap.
\end{enumerate}
\end{frame}

% =========================================================
% RANGKUMAN
% =========================================================
\begin{frame}{Rangkuman}
\small
\begin{itemize}
    \item \asm{MOV} menyalin tanpa menghapus sumber; legalitasnya dibatasi
          ukuran operand yang sama, larangan memori ke memori, dan
          larangan nilai segera ke register segmen.

    \vspace{0.2em}
    \item Pemindahan yang ditolak diatasi dengan register perantara.
          \asm{MOV} memindahkan isi memori, \asm{LEA} memindahkan
          alamatnya; \asm{XCHG} menukar isi dua operand.

    \vspace{0.2em}
    \item \asm{PUSH} dan \asm{POP} bekerja pada satuan word, menolak nilai
          segera, dan harus berpasangan seimbang agar \reg{SP} kembali.

    \vspace{0.2em}
    \item Instruksi aritmatika dasar meninggalkan jejak pada \eng{flag},
          kecuali \flag{CF} yang tidak disentuh \asm{INC} dan \asm{DEC}.
          \flag{CF} untuk operasi tak bertanda, \flag{OF} untuk bertanda.

    \vspace{0.2em}
    \item \asm{MUL}/\asm{IMUL} memakai operand implisit \reg{AL} atau
          \reg{AX} dan menghasilkan dua kali ukuran operand;
          \asm{DIV}/\asm{IDIV} menuntut pemeriksaan \asm{CMP AH, BL}.

    \vspace{0.2em}
    \item Aritmatika bertingkat disusun dengan \asm{ADC} dan \asm{SBB};
          bilangan biner diubah menjadi digit desimal dengan pembagian
          berulang sepuluh.
\end{itemize}
\end{frame}

% ---------------------------------------------------------
\begin{frame}{Cek Pemahaman}
\small
\begin{enumerate}
    \item Instruksi berikut yang \emph{tidak} sah pada prosesor 8086
          adalah \dots
          \begin{enumerate}\footnotesize
              \item \asm{MOV AX, BX} \quad
              \item \asm{MOV AX, nilai1} \quad
              \item \asm{MOV nilai1, nilai2} \quad
              \item \asm{MOV nilai1, AX}
          \end{enumerate}

    \vspace{0.25em}
    \item Instruksi yang \emph{tidak} mengubah nilai \flag{CF} adalah
          \dots
          \begin{enumerate}\footnotesize
              \item \asm{ADD AX, BX} \quad \item \asm{SUB AX, BX} \quad
              \item \asm{INC AX} \quad \item \asm{NEG AX}
          \end{enumerate}

    \vspace{0.25em}
    \item Kesalahan \eng{divide overflow} pada \asm{DIV BL} dapat dicegah
          dengan memeriksa \dots
          \begin{enumerate}\footnotesize
              \item apakah \reg{AL} bernilai nol
              \item apakah \reg{AH} lebih kecil daripada \reg{BL}
              \item apakah \flag{ZF} bernilai satu
              \item apakah \reg{BL} lebih kecil daripada \reg{AL}
          \end{enumerate}
\end{enumerate}
\end{frame}

% =========================================================
% PENUTUP
% =========================================================
\begin{frame}{Penutup}
\small
\begin{block}{Sebelum pertemuan berikutnya}
\begin{itemize}
    \item Selesaikan Tugas~4: \file{kalkulator.asm} beserta tangkapan
          layar keempat operasi dan pesan pembagi nol.
    \item Lengkapi tabel pengamatan \reg{AX}, \flag{CF}, \flag{ZF}, dan
          \flag{OF} serta jurnal praktikum.
    \item Ulangi bagian yang belum dikuasai dengan menelusuri
          \file{kali-bagi.asm} satu instruksi demi satu.
    \item Baca \textbf{Bab XI} --- Operasi Logika, Pergeseran, dan
          Rotasi.
\end{itemize}
\end{block}

\vspace{0.3em}
{\footnotesize\textbf{Bacaan lanjutan:} Bab X bagian Operand Implisit,
Makna Flag Hasil Aritmatika, dan Kesalahan Divide Overflow; aturan
\eng{flag} rinci pada dokumentasi Intel 8086; pembahasan interupsi dan
layanan DOS pada Bab XII.\par}

\vspace{0.35em}
\begin{center}
{\color{obeBlue}\sffamily\large\bfseries Terima kasih\par}
\vspace{0.15em}
{\color{obeGray}\sffamily\small \NamaMK\ (\KodeMK) \textbullet\
\SemesterTeks\par}
\end{center}
\end{frame}

\end{document}
